% ============================================================================
% Fichier  : partie3/P03_C08_normes_incident_preuve.tex
% Titre    : Normes de la Gestion d'Incident et de la Preuve
% Auteur   : MINKA MI NGUIDJOÏ Thierry Emmanuel
% ============================================================================
% Sources : ISO/IEC 27035 ; ISO/IEC 27037 ; NIST SP 800-86 ; RFC 3227.
% ============================================================================

\chapter{Normes de la Gestion d'Incident et de la Preuve}
\label{chap:normes}
\chaptermark{Normes de l'incident et de la preuve}

\epigraph{%
    « Une norme ne remplace pas le jugement~;\\
    elle en garantit la méthode et la trace. »%
}{Principe du cours}

\niveauChapitre{Fondamental}{Partie~II~; notions de gestion de la sécurité.}

\begin{boxobjectifs}
\begin{itemize}
    \item Décrire le cycle de gestion d'incident de l'ISO/IEC 27035.
    \item Connaître les principes de collecte de preuve de l'ISO/IEC 27037.
    \item Situer le guide NIST SP 800-86 et la RFC 3227.
    \item Articuler ces référentiels entre eux.
    \item Relier chaque norme à une pratique du cours.
\end{itemize}
\end{boxobjectifs}

\bigskip

% ============================================================================
\section*{Question directrice : quelles normes encadrent ce que LiZbiZ doit faire, du constat à la preuve~?}
% ============================================================================

Face à l'incident, LiZbiZ et NYONG SÉCURITÉ ne partent pas d'une page
blanche~: des normes internationales fixent le cycle de gestion et les
exigences de la preuve. Ce chapitre les présente et montre comment elles
s'articulent.

\medskip
\textbf{Plan du chapitre.}
\begin{enumerate}
    \item Le cycle ISO/IEC 27035 (section~\ref{sec:c08-27035}).
    \item La preuve selon l'ISO/IEC 27037 (section~\ref{sec:c08-27037}).
    \item NIST SP 800-86 et RFC 3227 (section~\ref{sec:c08-nist}).
    \item Articuler les référentiels (section~\ref{sec:c08-articuler}).
\end{enumerate}

\separateur

% ============================================================================
\section{Le Cycle ISO/IEC 27035}
\label{sec:c08-27035}
% ============================================================================

\begin{boxdef}
\termecle{ISO/IEC 27035}~: norme de gestion des incidents de sécurité de
l'information. Elle décrit un cycle en cinq phases~: \emph{planifier et
préparer}~; \emph{détecter et signaler}~; \emph{évaluer et décider}~;
\emph{répondre}~; \emph{tirer les leçons}.
\end{boxdef}

\begin{tableaucours}{L{4.0cm}L{8.2cm}}{%
    \tableauHeader{Phase (ISO 27035)} &
    \tableauHeader{Ce qu'elle recouvre}
}
Planifier et préparer & Politique, rôles, outils, playbooks \\
Détecter et signaler & Supervision, alertes, point de contact \\
Évaluer et décider & Qualification, gravité, décision de réponse \\
Répondre & Confinement, éradication, rétablissement \\
Tirer les leçons & Retour d'expérience, amélioration \\
\end{tableaucours}
\vspace{-0.8em}
\begin{center}{\small\itshape Tableau~8.1. Le cycle en cinq phases de l'ISO/IEC 27035~\cite{iso27035}}\end{center}

% ============================================================================
\section{La Preuve selon l'ISO/IEC 27037}
\label{sec:c08-27037}
% ============================================================================

\begin{boxdef}
\termecle{ISO/IEC 27037}~: lignes directrices pour l'identification, la
collecte, l'acquisition et la préservation des preuves numériques. Elle
repose sur trois exigences~: \termecle{pertinence}, \termecle{fiabilité} et
\termecle{suffisance}, garanties par la traçabilité des actes.
\end{boxdef}

\begin{boxprocedure}
\textbf{Préserver une preuve (ISO/IEC 27037)}
\begin{enumerate}
    \item Identifier le support et sa volatilité.
    \item Acquérir selon l'ordre de volatilité (du plus volatil au plus
          stable).
    \item Calculer un condensat (empreinte) à l'acquisition.
    \item Documenter chaque acte~: qui, quoi, quand, comment.
    \item Conserver en garantissant l'intégrité (chaîne de custody,
          Chapitre~\ref{chap:fondements}).
\end{enumerate}
\end{boxprocedure}

% ============================================================================
\section{NIST SP 800-86 et RFC 3227}
\label{sec:c08-nist}
% ============================================================================

Le \termecle{NIST SP 800-86} intègre la forensique à la réponse à incident en
quatre temps~: \emph{collecte}, \emph{examen}, \emph{analyse},
\emph{rapport}. La \termecle{RFC 3227} donne les lignes directrices de
collecte et d'archivage des preuves, dont le fameux \termecle{ordre de
volatilité} (mémoire avant disque, disque avant archives).

\begin{boxCRO}
\textbf{Analyse CRO de l'ordre de volatilité.} Capturer la mémoire d'abord
sert la \cro{0}{+}{+}~: on fige une trace fugace (fiabilité) exploitable
devant un tiers (opposabilité), sans exposer plus que nécessaire. Attendre,
c'est perdre la trace et son opposabilité.
\end{boxCRO}

% ============================================================================
\section{Articuler les Référentiels}
\label{sec:c08-articuler}
% ============================================================================

\begin{boxscenario}{Opération FARO. Quelle norme pour quel acte}
\textbf{Constatation.} Pour l'enquête, M.~BIKORO suit l'ISO 27035 (cycle),
l'ISO 27037 (préservation), le NIST 800-86 (démarche forensique) et la
RFC 3227 (ordre de volatilité).

\medskip
\textbf{Analyse} (confiance élevée). Les référentiels ne se concurrencent
pas~: 27035 cadre la gestion, 27037 et RFC 3227 encadrent la preuve, 800-86
relie forensique et réponse. Ensemble, ils garantissent une méthode
opposable.
\end{boxscenario}

\begin{boxmetier}{Analyste forensique et RSSI}
\textbf{Citer la norme suivie.} Dans un rapport, l'analyste indique quelle
norme a guidé chaque acte~: c'est ce qui rend la démarche vérifiable et
opposable. « J'ai acquis selon l'ISO 27037 » vaut mieux que « j'ai copié le
disque ».
\end{boxmetier}

\separateur

\begin{boxresume}
\textbf{Ce qu'il faut retenir de ce chapitre~:}
\begin{itemize}
    \item \textbf{ISO/IEC 27035}~: cycle d'incident en cinq phases.
    \item \textbf{ISO/IEC 27037}~: pertinence, fiabilité, suffisance de la
          preuve.
    \item \textbf{NIST 800-86}~: collecte, examen, analyse, rapport.
    \item \textbf{RFC 3227}~: ordre de volatilité.
    \item \textbf{Les référentiels s'articulent}~; on cite celui que l'on
          suit.
\end{itemize}
\end{boxresume}

\begin{boxinfo}
\textbf{Points de contrôle de l'élève.} Je sais~: énoncer les cinq phases de
27035~; citer les exigences de 27037~; décrire la démarche 800-86~; énoncer
l'ordre de volatilité~; articuler les référentiels.
\end{boxinfo}

\bigskip

\begin{boxexo}[title={\ding{45}\quad Exercice~8.1. Cinq phases \niveaubase}]
Citez les cinq phases du cycle ISO/IEC 27035 dans l'ordre.
\end{boxexo}

\begin{boxexo}[title={\ding{45}\quad Exercice~8.2. Exigences \niveaubase}]
Quelles sont les trois exigences de la preuve selon l'ISO/IEC 27037~?
\end{boxexo}

\begin{boxexo}[title={\ding{45}\quad Exercice~8.3. Volatilité \niveaubase}]
Dans quel ordre acquérir~: disque, mémoire, archives~?
\end{boxexo}

\begin{boxexo}[title={\ding{45}\quad Exercice~8.4. 800-86 \niveaumoyen}]
Décrivez les quatre temps du NIST SP 800-86 sur l'exemple de FARO.
\end{boxexo}

\begin{boxexo}[title={\ding{45}\quad Exercice~8.5. CRO \niveaumoyen}]
Expliquez l'analyse CRO de l'ordre de volatilité.
\end{boxexo}

\begin{boxexo}[title={\ding{45}\quad Exercice~8.6. Articulation \niveaumoyen}]
Quelle norme cadre la gestion~? Laquelle encadre la preuve~?
\end{boxexo}

\begin{boxexo}[title={\ding{45}\quad Exercice~8.7. Condensat \niveaumoyen}]
Pourquoi calculer un condensat dès l'acquisition~?
\end{boxexo}

\begin{boxexo}[title={\ding{45}\quad Exercice~8.8. Citer la norme \niveaumoyen}]
Pourquoi « acquis selon l'ISO 27037 » vaut-il mieux que « copié le disque »~?
\end{boxexo}

\begin{boxexo}[title={\ding{45}\quad Exercice~8.9. Démarche \niveauexpert}]
Rédigez, pour FARO, la démarche d'enquête en reliant chaque acte à sa norme.
\end{boxexo}

\begin{boxexo}[title={\ding{45}\quad Exercice~8.10. Playbook \niveauexpert}]
Intégrez l'ISO 27035 dans un playbook de réponse au rançongiciel de LiZbiZ.
\end{boxexo}

\bigskip

\begin{boxinfo}
\textbf{Pour aller plus loin.}
\begin{itemize}
    \item ISO/IEC 27035~\cite{iso27035}~; ISO/IEC 27037~\cite{iso27037}~;
          NIST SP~800-86~\cite{nist80086}~; RFC~3227~\cite{rfc3227}.
\end{itemize}
\end{boxinfo}

\begin{boxinfo}
\textbf{Transition.} Aux normes s'ajoute le droit~: le
Chapitre~\ref{chap:droit} traite du cadre juridique camerounais et africain.
\end{boxinfo}
